% mode: ini-etex
% Aimed at find_font_dimen, every branch: n<=0, within the font, past the end of the last font (it grows), past the end of
% an earlier font (the error), in reading and in writing; the reset of font_glue when 2<=n<=4 is written; and
% scan_font_ident's forms: \font, a font name, \textfont and the like, and "Missing font identifier".
\input prelude
\errorcontextlines=100
\font\a=cmr10 at 8pt \font\b=cmbx10 at 11pt
\def\say#1{\message{<#1>}}
% Reading.
\say{\the\fontdimen1\a}
\say{\the\fontdimen7\a}
\say{\the\fontdimen8\a}
\say{\the\fontdimen0\a}
\say{\the\fontdimen-1\a}
\say{\the\fontdimen-2147483647\a}
\say{\the\fontdimen2147483647\a}
\say{\the\fontdimen7\b}
\say{\the\fontdimen8\b}
% The last font grows, on reading too, and an earlier one does not.
\say{\the\fontdimen9\b}
\say{\the\fontdimen8\b}
\say{\the\fontdimen20\b}
\say{\the\fontdimen8\a}
\fontdimen8\a=1pt
\fontdimen0\a=1pt
\fontdimen-5\a=1pt
\fontdimen21\b=2pt
\say{\the\fontdimen21\b}
\say{\the\fontdimen22\b}
% Nullfont has no parameters of its own, and it is not the last font.
\say{\the\fontdimen1\nullfont}
\fontdimen1\nullfont=1pt
\font\c=cmr10 at 12pt
\say{\the\fontdimen8\c}
\fontdimen30\c=3pt
\font\d=cmr10 at 13pt
\fontdimen31\c=3pt
\say{\the\fontdimen30\c}
\say{\the\fontdimen31\c}
\say{\the\fontdimen8\d}
% The space glue is kept per font in font_glue: a change of 2, 3 or 4 throws it away, and others do not.
\def\spaces#1{\setbox1\hbox{#1a b}\showbox1 }
\tenrm \spaces{}
\fontdimen2\tenrm=4pt \spaces{}
\fontdimen3\tenrm=5pt \spaces{}
\fontdimen4\tenrm=6pt \spaces{}
\fontdimen5\tenrm=7pt \spaces{}
\fontdimen6\tenrm=8pt \spaces{}
\fontdimen7\tenrm=9pt \spaces{}
\font\e=cmr10 at 9pt
\fontdimen2\e=1pt \fontdimen1\e=1pt \spaces\e
\tenrm
% scan_font_ident.
\say{\the\fontdimen1\font}
\say{\the\fontdimen1\tenrm}
\say{\the\fontdimen1\textfont2}
\say{\the\fontdimen1\scriptfont1}
\say{\the\fontdimen1\scriptscriptfont0}
\say{\the\fontdimen1\textfont16}
\say{\the\fontdimen1\textfont15}
\say{\the\fontdimen1 x}
\say{\the\fontdimen1 \relax}
\say{\the\fontdimen1 \undefinedname}
\say{\the\fontdimen1 1}
\say{\the\fontdimen1 {}}
\def\f{\tenrm}\say{\the\fontdimen1 \f}
\def\f{  \tenrm}\say{\the\fontdimen1 \f}
\say{\the\fontdimen1    \tenrm}
\say{\the\hyphenchar\tenrm} \hyphenchar\tenrm=127 \say{\the\hyphenchar\tenrm}
\say{\the\skewchar\tenrm} \skewchar\tenrm=-5 \say{\the\skewchar\tenrm}
\hyphenchar x=1
\skewchar\relax=1
\font\l=cmr10 at 7pt
\letterspacefont\m\l 100 \say{\the\fontdimen1\m}
\pdfcopyfont\n\l \say{\the\fontdimen1\n}
\let\o\font \say{\the\fontdimen1\o}
\end
